\section{Exact homology in a moving degree window}
\label{sec:window-theory}

The next optimization changes which chain groups are constructed.  It is
useful when the decision procedure needs a prescribed part of Khovanov
homology, and admits an unconditional complexity analysis for the delivered
implementation.  It does not assert that a bounded part of homology always
detects a nontrivial knot.

Throughout this section the coefficient field is $\mathbb F_2$.  The quantum
grading is forgotten, but the raw cube degree is retained.  A crossing
complex is supported in raw degrees $0,1$; delooping changes no homological
degree.  If $\nu$ denotes the number of negative crossings in
\texttt{Diagram.signs()}, normalized degree $h$ corresponds to raw degree
$h+\nu$.  This convention must be read from the oriented PD diagram:
\texttt{Diagram.from\_braid(2,[1])} has $\nu=1$ and raw homology in degree
$1$, while \texttt{Diagram.from\_braid(2,[-1])} has $\nu=0$ and raw homology
in degree $0$.  Both closures are unknots.  Counting negative integers in a
braid word instead would reverse this convention.

\subsection{The support lemma and its guard degrees}

For a cochain complex $C$, write $C[L,U]$ for its brutal interval truncation:
retain $C^j$ for $L\leq j\leq U$, retain the differentials between these
groups, and set every other group and differential to zero.  This is a
complex, but its homology at $L$ and $U$ need not equal that of $C$.

\begin{lemma}[Support under a bounded continuation]
\label{lem:bounded-continuation}
Suppose a remaining continuation is a composite of additive planar
extension operations whose total raw homological support is contained in
$[u,v]$, followed by closure and the Khovanov functor.  To preserve final
homology in raw degrees $[a,b]$, it suffices to replace the current complex
$C$ by
\[
 C[a-v-1,b-u+1].
\]
This replacement gives exactly the same final chain groups and
differentials in the three-degree neighborhood $[a-1,b+1]$.
\end{lemma}

\begin{proof}
Write the total continuation formally as $Q$.  A summand contributing in
final degree $j$ comes from current degree $p$ and continuation degree $q$
with $p+q=j$ and $u\leq q\leq v$.  If $a-1\leq j\leq b+1$, then
\[
 a-v-1\leq p\leq b-u+1.
\]
Consequently no removed current summand contributes to any of these final
chain groups.  The total differential is the sum of the current and
continuation differentials, with the usual tensor signs.  Its entries
between retained final groups are therefore unchanged as well.  Applying
the additive closure and Khovanov functors preserves these equalities.
Finally,
$H^j=\ker(d^j)/\operatorname{im}(d^{j-1})$ depends only on the groups in
degrees $j-1,j,j+1$ and the two indicated maps.
\end{proof}

Both guards are necessary in the general algebraic statement.  The
two-term complex $\mathbb F_2\xrightarrow{1}\mathbb F_2$ in degrees
$a-1,a$ has zero $H^a$, whereas discarding degree $a-1$ creates a false
class.  Placing the same complex in degrees $b,b+1$ shows the corresponding
problem at the upper end.  Keeping only the desired degrees is therefore
unsound even after the knot has closed.

For $n-i$ unprocessed crossings the continuation support is $[0,n-i]$.
Define
\begin{equation}
 I_i=[L_i,U_i]\cap[0,i],\qquad
 L_i=a-(n-i)-1,\qquad U_i=b+1.
 \label{eq:moving-window}
\end{equation}
The upper endpoint is fixed.  The lower endpoint moves upward by one after
every crossing, because progressively fewer future $1$-resolutions remain
available to reach the target.

\begin{theorem}[Repeated windowing with arbitrary unit cancellation]
\label{thm:window-correctness}
Start from the empty partial diagram, and process crossings in any valid
scan order.  After each extension, retain only the summands in $I_i$ and
the differential entries whose two endpoints are retained.  Delooping and
any sequence of valid Gaussian unit cancellations may be performed at each
stage.  The resulting closed complex has the exact original Khovanov
homology in every raw degree $a\leq j\leq b$.

The discarded summands may be excluded before allocation, provided the
constructed objects and maps equal the brutal truncation of the full next
extension.
\end{theorem}

\begin{proof}
Compare each intermediate complex only after applying its actual remaining
continuation.  At the beginning this gives the ordinary complete Khovanov
complex.  One crossing extension merely factors that continuation into its
first operation and its suffix, so it leaves the completed complex
unchanged.  Lemma~\ref{lem:bounded-continuation} shows that clipping to
$I_i$ preserves the completed homology in $[a,b]$.

A Gaussian cancellation is a chain homotopy equivalence of the current
complex.  Every remaining planar extension is additive and the induced
operation on total complexes preserves the equations defining a chain
homotopy.  Closure and the Khovanov functor do the same.  Hence a cancellation
also preserves the completed target homology.  Induction proves the claim.

This argument does not require the artificial endpoint homology of a
truncation to be the original endpoint homology.  Its possible effects are
already outside the final target and its adjacent differential groups.
Nor does it require the same pivots as an untruncated run.  Pre-allocation
pruning realizes the same truncated complex by omitting objects and maps
that the support lemma would immediately delete.
\end{proof}

\paragraph{Operational form.}
The following sequence specifies the mathematical operation independently
of the pivot heuristic:
\begin{enumerate}
\item Fix the full validated crossing order and target raw interval $[a,b]$.
\item At step $i$, compute the band $I_i$ using the original total $n$.
\item For each current object of degree $h$ and each smoothing
  $\epsilon\in\{0,1\}$, allocate its delooped children only if
  $h+\epsilon\in I_i$.
\item Construct the crossing saddles and transferred old differentials only
  when both endpoint children were retained.
\item Cancel invertible entries, then continue to the next crossing.
\item At closure compute homology using both adjacent differentials, and
  return only the degrees in $[a,b]$.
\end{enumerate}
The frontier itself belongs to the processed subdiagram.  It must advance
even when the retained algebraic complex is empty.  A resource cap is a cap
on the retained allocation; exhaustion gives an incomplete computation,
never a homology value.

\begin{lemma}[Removing isolated guards]
\label{lem:isolated-window-guards}
After reduction at a prefix with $r$ crossings remaining, an isolated
object of raw degree $h$ may be removed whenever
$[h,h+r]\cap[a,b]=\varnothing$.  A nonisolated guard cannot be removed by
this test.
\end{lemma}

\begin{proof}
An object with no incoming or outgoing nonzero differential entry is an
independent direct summand.  Completing this summand through the remaining
crossings gives a complex supported in $[h,h+r]$.  It therefore contributes
zero homology to the target band.  Direct sums commute with the
continuation, so deleting the summand is valid.  For a nonisolated object
an incident differential may create or kill a target class, as the
two-term examples above demonstrate.
\end{proof}

With the one-degree guards in~\eqref{eq:moving-window}, a differential
component whose entire degree support misses $[a-r,b]$ lies in one guard
degree.  Since differentials change raw degree by one, such components
are already isolated objects.  Thus this particular additional pruning
does not require a general graph search.

\subsection{An unconditional bound for the implemented scanner}

Let $s_0(D)$ and $s_1(D)$ denote the numbers of circles in the complete
all-$0$ and all-$1$ smoothings of a nonempty diagram $D$.  These quantities
are computed by pairing the appropriate crossing slots and counting
components.  Let $w$ be the largest number of frontier endpoints in the
chosen order.  No assumption that these endpoints lie on one boundary
circle is made.

For $0\leq k\leq n$, put
\begin{equation}
 Q(n,k)=\sum_{j=0}^{k}2^j\binom{n}{j},\qquad
 B_0(D;a,b)=2^{s_0(D)}Q\bigl(n,\min\{n,b+1\}\bigr).
 \label{eq:elementary-window-bound}
\end{equation}
Targets are first intersected with $[0,n]$; an empty intersection has zero
homology.  The separate input convention that an empty PD represents one
circle is handled directly by the implementation.

\begin{lemma}[Circle and object counts]
\label{lem:window-circle-count}
Every raw cube state of a prefix with $j$ $1$-resolutions has at most
$s_0(D)+j$ already closed circles.  At every stage the number of retained
physical summands in the window scanner, including the allocation before
Gaussian reduction, is at most $B_0(D;a,b)$.
\end{lemma}

\begin{proof}
Complete the prefix state with $0$-resolutions at the remaining crossings.
Every circle already closed in the prefix remains a distinct circle of
that completion.  The completed state differs from the all-$0$ state at
$j$ crossings.  Changing one smoothing is a surgery on two arcs and changes
the total circle count by at most one.  The completed state therefore has
at most $s_0(D)+j$ circles.

A state with $c$ already closed circles contributes $2^c$ summands after
delooping.  There are $\binom{i}{j}$ states of degree $j$ at prefix length
$i$.  Extension and delooping enumerate the usual cube objects; Gaussian
cancellation and truncation only remove objects.  More formally, surviving
objects can be assigned distinct original cube-basis labels, and extending
these labels and adding new circle labels preserves this assignment.
Consequently the degree-$j$ allocation has size at most
$2^{s_0(D)+j}\binom{i}{j}$.  Every retained degree satisfies
$j\leq b+1$.  Summing and using $\binom{i}{j}\leq\binom{n}{j}$ proves the
bound.
\end{proof}

One can keep a more informative prefix bound.  If $z_i$ is the number of
closed circles in the all-$0$ smoothing of the processed prefix, then the
same local surgery argument gives
\begin{equation}
 \#\text{objects after allocation at }i
 \leq\sum_{j\in I_i}2^{z_i+j}\binom{i}{j}.
 \label{eq:prefix-window-bound}
\end{equation}
This explicitly exhibits the saving from the moving lower boundary.
Since $z_i\leq s_0(D)$, it implies the uniform bound above.

\begin{theorem}[Bit complexity with unrestricted scan frontiers]
\label{thm:window-cost}
For the delivered $\mathbb F_2$ dotted-cobordism representation there is an
absolute constant $c$ such that the uncapped window computation takes at
most
\begin{equation}
 2^{c(w+1)}(n+1)^c\bigl(B_0(D;a,b)+1\bigr)^c
 \label{eq:window-cost-exact}
\end{equation}
bit operations, after the normalized PD diagram and order are supplied.
In particular, if $k=\min\{n,b+1\}$, its time is
\begin{equation}
 2^{O(w+s_0(D)+k)}(n+1)^{O(k+1)}.
 \label{eq:window-cost-asymptotic}
\end{equation}
The implied constants do not depend on the knot, the order, or $k$.
\end{theorem}

\begin{proof}
Two matchings on at most $w$ endpoints have at most $w/2$ circle components
in their union.  The implemented dot-mask basis of the corresponding
morphism space therefore has at most $2^{w/2}$ elements.  A morphism is
stored as at most that many bits.  The compiled gluing plans use
polynomially many boundary records; enumerating their monomials, composing
maps and evaluating neck-cutting sums takes $2^{O(w)}$ bit operations.
Adding a crossing introduces only four local slots, which changes this
bound by a constant factor in the exponent.

With $B=B_0(D;a,b)$ objects, there are at most $B^2$ differential entries.
At most $B/2$ unit cancellations occur in a stage; each cancellation
examines at most $B^2$ pairs of incident entries.  Bookkeeping, the lazy
pivot heap, stage-local caches and equality checks add only polynomial
factors in $n,B,w$.  The implementation uses $\mathbb F_2$, so coefficient
growth contributes no additional integer-height parameter.  There are
$n$ stages.  Lemma~\ref{lem:window-circle-count} now gives
\eqref{eq:window-cost-exact}.  Finally,
\[
 Q(n,k)\leq(k+1)2^k(n+1)^k,
\]
which yields~\eqref{eq:window-cost-asymptotic}.

This count never bounds the number of all possible matchings by a Catalan
number.  It bounds the matchings actually attached to retained cube
summands, and then bounds the algebra attached to each pair.  Thus the
frontier issue for arbitrary crossing orders does not enter the proof.
\end{proof}

\subsection{Choosing the less expensive mirror}

Mirroring swaps the two smoothings.  Over a field, the closed Khovanov
complex of the mirror is dual to the original complex, with raw degrees
reversed:
\[
 H^j_{\mathrm{raw}}(D^{!};\mathbb F_2)
 \cong H^{n-j}_{\mathrm{raw}}(D;\mathbb F_2)^*.
\]
The underlying crossing graph, and hence the frontier size in a fixed
crossing order, is unchanged.  Also $s_0(D^{!})=s_1(D)$.  Define
\begin{equation}
 B_1(D;a,b)=2^{s_1(D)}Q\bigl(n,\min\{n,n-a+1\}\bigr).
 \label{eq:mirror-window-bound}
\end{equation}

\begin{corollary}[A computable choice of direction]
\label{cor:mirror-window}
Exact raw homology in $[a,b]$ can be computed with the bound
\eqref{eq:window-cost-exact} after replacing $B_0$ by
$\min\{B_0,B_1\}$.  If $B_1<B_0$, run on the mirror with target
$[n-b,n-a]$, then reverse the returned degree labels.
\end{corollary}

\begin{proof}
Apply Theorem~\ref{thm:window-cost} to either $D$ or $D^{!}$, using the
displayed duality.  The all-state circle counts are found by finite
component calculations.  Both $Q$ values have $O(n)$ bits, since
$Q(n,k)\leq3^n$, so comparing the two explicit bounds has polynomial cost.
\end{proof}

This choice accounts for the actual circle counts.  Merely choosing the
smaller extreme depth can miss a large difference between $s_0$ and $s_1$.
The choice minimizes this proved upper bound; it does not promise to
predict the faster observed run on every input.

Let $s$ be the number of Seifert circles, and let
$\mu=\min\{\nu,n-\nu\}$.  The oriented smoothing differs from the all-$0$
smoothing at exactly $\nu$ crossings and from the all-$1$ smoothing at
exactly $n-\nu$ crossings in the stated PD convention.  The same one-surgery
circle bound gives
\[
 s_0\leq s+\nu,\qquad s_1\leq s+(n-\nu).
\]
For normalized target degrees $[-d,d]$, one may therefore choose a mirror
and obtain the following useful uniform consequence.

\begin{corollary}[Minority-sign and Seifert-circle parameters]
\label{cor:seifert-window-cost}
For $d\geq0$, exact Khovanov homology in normalized degrees $[-d,d]$ can be
computed by the window scanner in
\begin{equation}
 2^{O(w+s+\mu+d)}(n+1)^{O(\mu+d+1)}
 \label{eq:seifert-window-cost}
\end{equation}
bit operations.  In particular this is quasi-polynomial when
$w+s=O(\log^2(n+2))$ and $\mu+d=O(\log(n+2))$.

For a knot given as the closure of a $t$-strand braid scanned in braid
order, $s=t$ and $w\leq2t$.  The bound becomes
\[
 2^{O(t+\mu+d)}(n+1)^{O(\mu+d+1)}.
\]
It is polynomial for fixed $t,\mu,d$, and quasi-polynomial for
$t=O(\log^2(n+2))$ and $\mu+d=O(\log(n+2))$.
\end{corollary}

\begin{proof}
On the side with $\mu$ negative PD crossings, the requested raw upper
degree is at most $\mu+d$, and its upper guard is at most $\mu+d+1$.
Insert $s_0\leq s+\mu$ in Theorem~\ref{thm:window-cost}.  The logarithm of
the resulting time is
$O(w+s+\mu+d+(\mu+d+1)\log(n+2))$, proving the stated regimes.

Oriented smoothing of a braid closure yields its $t$ Seifert circles.  In
braid order the processed prefix has at most one top and one current end
per strand, so its frontier has at most $2t$ endpoints.  Substitution gives
the braid assertions.  The sign convention for braid-word generators does
not affect $\mu$, although it does affect which mirror is selected.
\end{proof}

\subsection{Scope and comparison with existing extremal algorithms}

An exact band can reject an unknot whenever it differs from the unknot's
homology in that band.  In particular, excess dimension in normalized
degree zero, or nonzero homology in a computed nonzero degree, is a sound
certificate of nontriviality.  An agreeing band remains inconclusive.
For example, the trefoil has the unknot's dimension in degree zero and has
additional homology elsewhere.  A complete recognizer must use a fallback,
eventually cover every degree, or prove a separate support or witness
theorem for the input class.

Kelom\"aki--Sch\"utz give an extremal-degree algorithm based on special
``nice'' scanning sequences and local cancellations
\cite{kelomaki}.  Their Lemma~5.5 bounds intermediate degree-$j$ summands by
$\binom{i}{j}$, with a factor two at closure; their finite-field
Algorithm~3 permits exhaustive Gaussian cancellation.  This can remove
the all-state circle factor in an appropriate geometric assembly.  It is
a distinct guarantee from the elementary object bound proved here.  The
current implementation accepts arbitrary valid crossing orders, and no
claim that every such order satisfies their nice-scan hypotheses is made.

The cited work also proves polynomial-time computation of full homology
for three-braid closures by specialized structure, despite exponential
explicit-basis examples.  Neither that specialized engine nor its general
fixed-strand extension is supplied by degree clipping.  Our contribution
here is an exact moving interval implemented before allocation, together
with a uniform, independently proved bound for that implementation.

\paragraph{Relation to \texttt{openai/math}.}
The source of \emph{Parabolic intersections in Artin groups} in
\texttt{openai/math} uses the principle that later functors
with bounded level amplitude cannot bring discarded layers back to an
extreme level; see its lemma ``Composition at the bottom'' in
\texttt{build/sections/03\_layers.tex} \cite{openai_artin}.  Its categorical-action section
also retains explicit chain homotopies before tensoring.  These are useful
methodological parallels.  The Artin-module levels are not the raw
Khovanov degrees used here, and the present proof imports no theorem about
Artin groups or their parabolic subgroups.  Lemma~\ref{lem:bounded-continuation}
supplies the required statement directly in the scanner's category.  The
reviewed repository revision is
\texttt{adc7f1241b42e322a6451854ab7e4b4c146bf78a}.
